\documentclass[a4paper,11pt]{article}
\usepackage[ngerman]{babel}
\usepackage[utf8]{inputenc}
\usepackage[T1]{fontenc}
\usepackage{amsmath,amssymb}
\usepackage{booktabs}
\usepackage[margin=2.2cm]{geometry}
\usepackage{url}

\title{Peridynamische Modelle im 1D-Stab -- Kurzfassung}
\author{BBVV1D.jl}
\date{\today}

\begin{document}
\maketitle

Aufbau, Bezeichnungen und Zeitintegration folgen
Peridynamics.jl\footnote{\url{https://github.com/kaipartmann/Peridynamics.jl}}; die 1D-Formeln sind die dortigen 3D-Formeln
mit $E$ statt $3K$ bzw.\ $\alpha$ statt $15G$.

\section*{Notation}
Punkte $X_i$ mit Volumen $V_j=\Delta x$, Horizont $\delta$ (hier $\delta = 3.015\,\Delta x$), Bonds $j\in\mathcal H_i$ mit
\[
  \xi_{ij} = X_j - X_i,\qquad L_{ij} = |\xi_{ij}|,\qquad l_{ij} = |x_j - x_i|,\qquad
  e_{ij} = l_{ij} - L_{ij},\qquad \varepsilon_{ij} = e_{ij}/L_{ij}.
\]
Kernel (Influence-Funktion) $\omega(L)$: konstant $1$, linear $\delta/L$ (Standard wie in Peridynamics.jl)
oder kubischer B-Spline. Momente der Bondlängen pro Punkt:
\[
  M_{n,i} = \sum_{j\in\mathcal H_i} L_{ij}^n\,V_j,\qquad q_i = M_{1,i},\qquad
  m_i = \sum_{j} \omega_{ij} L_{ij}^2\,V_j .
\]
Im Inneren gilt bei $\delta = 3.015\,\Delta x$:

\begin{center}
\begin{tabular}{@{}lll@{}}
  \toprule
  & diskret & Kontinuum \\
  \midrule
  $q$   & $12\,\Delta x^2$ & $\delta^2 \approx 9.09\,\Delta x^2$ \\
  $M_2$ & $28\,\Delta x^3$ & $2\delta^3/3 \approx 18.3\,\Delta x^3$ \\
  $M_3$ & $72\,\Delta x^4$ & $\delta^4/2 \approx 41.3\,\Delta x^4$ \\
  \bottomrule
\end{tabular}
\end{center}

\paragraph{Bondkonstante und Zeitschritt.} $b_c = 2E/\delta^2$ (1D-Gegenstück zu $18K/(\pi\delta^4)$,
aus $\tfrac12 E\varepsilon^2 = \tfrac14 b_c\varepsilon^2\delta^2$). Für alle Modelle
$\Delta t = 0.7\,\min_i \sqrt{2\rho \big/ \textstyle\sum_j V_j\, b_c/L_{ij}}$ wie in Peridynamics.jl.

\section*{Modelle}

\paragraph{BB -- bond-based.}
\[
  b_i = \sum_{j} \beta_{ij}\, b_c\, \varepsilon_{ij}\,\frac{x_j-x_i}{l_{ij}}\,V_j .
\]
Ohne Korrektur ist $\beta = 1$; mit den Kontinuumskonstanten ist die Welle dann um $+15\,\%$ zu schnell.
\emph{Energiekorrektur} (Le \& Bobaru 2018, wie Peridynamics.jl):
$\beta_{ij} = \tfrac12(\mathrm{m}_i + \mathrm{m}_j)$ mit $\mathrm{m}_i = \Psi_\mathrm{theo}/\Psi_{\mathrm{PD},i}$
bei gleichmäßiger Dehnung. In 1D ergibt das genau
$\beta_{ij} b_c = E\,(1/q_i + 1/q_j)$, also die diskrete Kalibrierung inklusive Rand.

\paragraph{CC -- Energieäquivalenz (Mittelung über $\varepsilon^2$).}
\[
  \bar\varepsilon_i = c \sum_j L_{ij}\,\Delta u_{ij}\,V_j ,\qquad
  b_i = E\,c \sum_j \bigl(\bar\varepsilon_i+\bar\varepsilon_j\bigr) L_{ij}\,V_j ,
\]
\[
  c = \frac{3}{2\delta^3}\ \text{(Kontinuum, Original)}\qquad\text{oder}\qquad
  c = \frac{1}{\sqrt{q\,M_3}}\ \text{(diskret)}.
\]
Wegen der geraden Gewichtung $L_{ij}$ ist $\bar\varepsilon$ kein Dehnungsmaß: Für $u = \varepsilon X$ ist
$\bar\varepsilon = 0$ im Inneren, für quadratisches $u$ gilt $\bar\varepsilon = c\,M_3\,u''/2$.
Mit $c = 3/(2\delta^3)$ ergibt sich in 1D $b = 2.59\,E u''$ ($+61\,\%$ Wellengeschwindigkeit).
Die Energiekorrektur ist nicht anwendbar, weil $\Psi_\mathrm{PD} = 0$ bei gleichmäßiger Dehnung.

\paragraph{OSB -- ordinary state-based (LPS).}
\[
  \theta_i = \frac{1}{m_i}\sum_j \omega_{ij}\, L_{ij}\, e_{ij}\,V_j ,\qquad
  \Psi_i = \frac{E}{2}\theta_i^2 + \frac{\alpha}{2m_i}\sum_j \omega_{ij}\bigl(e_{ij} - \theta_i L_{ij}\bigr)^2 V_j ,
\]
\[
  t_{ij} = \omega_{ij}\,\bigl(c_2 L_{ij} + c_1 e_{ij}\bigr),\qquad c_1 = \frac{\alpha}{m_i},\quad
  c_2 = \theta_i\Bigl(\frac{E}{m_i} - c_1\Bigr),
\]
$t_{ij}$ wirkt auf $i$ mit $V_j$ und als Reaktion auf $j$ mit $-V_i$. Standard $\alpha = E$.
Da $\theta$ auch am Rand exakt ist, ist der Energie-Korrekturfaktor in 1D gleich $1$.

\paragraph{Äquivalenz.} Mit linearem Kernel und $\alpha = E$ fällt $\theta$ heraus:
$t_{ij} = E\,\delta\,\varepsilon_{ij}/m_i$ und $m_i = \delta\,q_i$. \textbf{OSB (linear, $\alpha=E$) ist in 1D
identisch mit BB mit Energiekorrektur} -- auch am Rand und in allen Rechnungen bis auf Rundung.

\section*{Dispersionsrelationen}
Mit $u_j = \hat u\,e^{\mathrm i(kX_j-\omega t)}$ ergibt sich im Inneren $\rho\,\omega^2 = \Phi(k)$ mit
\[
  S = \sum_j \omega_j \xi_j \sin(k\xi_j)\,V_j,\quad C = \sum_j \omega_j\bigl(1-\cos k\xi_j\bigr)V_j ,\quad
  P_\pm = \sum_j \bigl(1\pm\cos k\xi_j\bigr)L_j\,V_j ,
\]
\begin{align*}
  \Phi_{\mathrm{BB}} &= \frac{2E}{q}\sum_j \frac{1-\cos k\xi_j}{L_j}\,V_j , \qquad
  \Phi_{\mathrm{CC}} = E\,c^2\,P_+\,P_- , \qquad
  \Phi_{\mathrm{OSB}} = E\,\frac{S^2}{m^2} + \frac{\alpha}{m}\Bigl(2C - \frac{S^2}{m}\Bigr) .
\end{align*}
Für $k\to 0$ gilt bei allen $\Phi \to E k^2$ (diskret kalibriert). Für lange Wellen
$\omega/(ck) \approx 1 - \beta\,(k\Delta x)^2$ mit
\[
  \beta_{\mathrm{BB}} = \frac{M_3}{24\,q},\qquad
  \beta_{\mathrm{CC}} = \frac{M_5}{24\,M_3} + \frac{M_3}{8\,q},\qquad
  \beta_{\mathrm{OSB}} = \frac{M^\omega_4}{M^\omega_2}\Bigl(\frac16 - \frac{\alpha}{8E}\Bigr),\quad M^\omega_n = \sum_j \omega_j L_j^n V_j .
\]
Für $\alpha = 4E/3$ verschwindet der $k^2$-Term bei OSB (für jeden Kernel).
CC hat für $\delta \ge 2\Delta x$ keine Nullenergie-Mode ($P_+ > 0$), OSB nicht für $\alpha > 0$.
Die Formeln sind numerisch gegen die Kraftfunktionen geprüft.

\section*{Vergleich}
\begin{center}
\small
\begin{tabular}{@{}lccccc@{}}
  \toprule
  & BB & CC & \multicolumn{3}{c}{OSB, Kernel $\omega$} \\
  \cmidrule(l){4-6}
  & Energiek. & diskret & konst. & linear & B-Spline \\
  \midrule
  $\beta$ (lange Wellen)          & $0.25$  & $1.07$ & $0.29$ & $0.25$ & $0.086$ \\
  Phasenfehler $<1\,\%$ ab $\lambda/\Delta x$ & $31$ & $65$ & $34$ & $31$ & $18$ \\
  Impulserhaltung                 & ja      & nein (Rand) & ja & ja & ja \\
  max.\ Fehler, $l/1000$\textsuperscript{*} & $0.037$ & $0.35$ & $0.045$ & $0.037$ & $0.012$ \\
  max.\ Fehler, $l/2000$\textsuperscript{*} & $0.010$ & $0.14$ & $0.012$ & $0.010$ & $0.006$ \\
  \bottomrule
\end{tabular}
\end{center}
{\footnotesize\textsuperscript{*}PMMA ($E = 3.2$\,GPa, $\rho = 1200$\,kg/m$^3$), Stab $1$\,m, Geschwindigkeitspuls über eine
Randschicht der Dicke $\delta$, Pulslänge $l/10$, freies Ende, Laufzeit $2l/c$;
Fehler $\max_t \Vert u-u_\mathrm{ref}\Vert/\max\Vert u_\mathrm{ref}\Vert$ gegen d'Alembert.
Ohne Korrektur (BB) bzw.\ mit $c = 3/(2\delta^3)$ (CC) liegt der Fehler bei etwa $1.5$: Die Pulse überlappen nicht mehr.
Bei CC stammt etwa $0.1$ aus der Dispersion im Inneren, der Rest aus der Reflexion am freien Ende.
OSB mit $\alpha = 10E$ ist mit dem Zeitschritt aus $b_c$ instabil.}

\paragraph{Abgleich mit Peridynamics.jl (3D).} Dünner Stab mit $500$ Punkten in Längsrichtung
($\Delta x = l/500$, $\delta = 3.015\,\Delta x$) und $4\times4$ bzw.\ $8\times8$ Punkten im Querschnitt, $\nu = 0.25$,
BB und OSB (linearer Kernel) jeweils mit Energiekorrektur, gleicher Puls über eine Randschicht der Dicke $\delta$,
$u_x$ über den Querschnitt gemittelt. Die Wellengeschwindigkeit ist aus der Laufzeit des Pulsmaximums zwischen
$x = 0.3\,l$ und $0.8\,l$ bestimmt, der Fehler wie oben gegen d'Alembert:

\begin{center}
\begin{tabular}{@{}lcccc@{}}
  \toprule
  & \multicolumn{2}{c}{$c/\sqrt{E/\rho}$} & \multicolumn{2}{c}{max.\ Fehler} \\
  \cmidrule(lr){2-3}\cmidrule(l){4-5}
  Querschnitt & BB & OSB & BB & OSB \\
  \midrule
  $4\times4$ ($8\,\mathrm{mm}$)  & $1.032$ & $1.048$ & $1.43$ & $1.44$ \\
  $8\times8$ ($16\,\mathrm{mm}$) & $0.999$ & $0.987$ & $0.24$ & $0.38$ \\
  1D, $\Delta x = l/500$         & $1.000$ & $1.000$ & $0.15$ & $0.15$ \\
  \bottomrule
\end{tabular}
\end{center}

Bei $4\times4$ Punkten ist der Querschnitt mit $8\,\mathrm{mm} \approx 1.3\,\delta$ schmaler als der Horizont
($6\,\mathrm{mm}$ Radius); die Welle ist um $3$--$5\,\%$ zu schnell, vermutlich weil die Querkontraktion teilweise behindert ist
(vollständig behindert: $\sqrt{M/E} = 1.095$ mit $M = E(1-\nu)/((1+\nu)(1-2\nu))$). Bei $8\times8$ Punkten
stimmt die Ankunftszeit, der Puls wird aber breiter und zieht Nachschwingungen hinter sich her, bei OSB stärker als bei BB.
Der verbleibende Fehler ist größer als in 1D bei gleichem $\Delta x$; mit dem Fehlermaß
lassen sich die Ankunftszeit und die Pulsverbreiterung nicht unterscheiden.
Nach diesen beiden Rechnungen sollte der Querschnitt für einen Vergleich mit 1D mindestens etwa $2.5\,\delta$ breit sein
($8\times8$: $16\,\mathrm{mm} \approx 2.7\,\delta$); feinere Querschnitte habe ich nicht gerechnet.

\paragraph{Randeffekt.} In allen Modellen gilt $b\neq 0$ unter homogener Dehnung in einem Streifen der Breite $2\delta$ am Rand,
weil die Momente der Nachbarn dort abweichen (peridynamischer Oberflächeneffekt).
Wird die Last nur an einem einzelnen Punkt aufgebracht, regt das kurzwellige Anteile an; über eine Randschicht der Dicke $\delta$ nicht.

\end{document}
